% from ega4-16-fr.tex
\oldpage[IV-4]{10}
compte de l'exactitude du foncteur $\rho^*$, on obtient un di-homomorphisme
d'Algèbres graduées (relatif à l'homomorphisme
$\lambda^\#:\rho^*(\mathcal{O}_Y)\to\mathcal{O}_{Y'}$)
\[
  \operatorname{gr}(u):
  \rho^*(\mathscr{G}r_\bullet(f))
  \longrightarrow\mathscr{G}r_\bullet(f')
  \tag{16.2.1.3}
  \label{IV.16.2.1.3}
\]
(ou, si on veut, un $\rho$-morphisme (0\textsubscript{I}, 3.5.1)
$\mathscr{G}r_\bullet(f)\to\mathscr{G}r_\bullet(f')$), et en particulier un
di-homomorphisme des faisceaux conormaux
\[
  \operatorname{gr}_1(u):
  \rho^*(\mathscr{G}r_1(f))\longrightarrow\mathscr{G}r_1(f').
\]
Il est immédiat en outre que ces homomorphismes donnent lieu à un diagramme
commutatif
\[
  \xymatrix{
    \rho^*(\mathscr{S}_{\mathcal{O}_Y}^{\bullet}
      (\mathscr{G}r_1(f)))
      \ar[r] \ar[d]_-{\mathscr{S}(\operatorname{gr}_1(u))}
      & \rho^*(\mathscr{G}r_\bullet(f))
        \ar[d]^-{\operatorname{gr}(u)} \\
    \mathscr{S}_{\mathcal{O}_{Y'}}^{\bullet}
      (\mathscr{G}r_1(f'))
      \ar[r]
      & \mathscr{G}r_\bullet(f')
  }
  \tag{16.2.1.4}
  \label{IV.16.2.1.4}
\]
où les flèches horizontales sont les homomorphismes canoniques (16.1.2.2).

Enfin, si l'on a un diagramme commutatif de morphismes d'espaces annelés
\[
  \xymatrix{
    Y \ar[r]^f & X \\
    Y' \ar[r]_{f'} \ar[u]^u & X' \ar[u]_v \\
    Y'' \ar[r]_{f''} \ar[u]^{u'} & X'' \ar[u]_{v'}
  }
\]
où $f''=(\psi'',\theta'')$ est tel que $\theta''^\#$ soit surjectif, et si
$w'_n$ et $w''_n$ sont définis à partir de $u'$, $v'$ d'une part, et de
$u''=u\circ u'$, $v''=v\circ v'$ de l'autre, alors on a
$w''_n=w_n\circ w'_n$, comme il résulte aussitôt des définitions et de
(0\textsubscript{I}, 3.5.5); on a de même
$\operatorname{gr}(u'')=
\operatorname{gr}(u')\circ\rho'^*(\operatorname{gr}(u))$ si
$u'=(\rho',\lambda')$. On peut donc dire que les $Y^{(n)}$ et
$\mathscr{G}r_\bullet(f)$ dépendent fonctoriellement de $f$.
\end{env}

\begin{proposition}[16.2.2]
\label{IV.16.2.2}
Avec les notations et hypothèses de (16.2.1), supposons de plus que $f$, $f'$,
$u$ et $v$ soient des morphismes de préschémas. Alors:
\begin{enumerate}
\item[(i)] Les morphismes $w_n:Y'^{(n)}\to Y^{(n)}$ sont des morphismes de
préschémas.
\item[(ii)] Si $Y'=Y\times_X X'$, $u$ et $f'$ étant les projections
canoniques, et si $f$ est une immersion ou si $v$ est plat, on a
$Y'^{(n)}=Y^{(n)}\times_X X'$.
\item[(iii)] Si $Y'=Y\times_X X'$ et si $v$ est plat (resp. si $f$ est une
immersion), l'homomorphisme
\[
  \operatorname{Gr}(u)=\operatorname{gr}(u)\otimes 1:
  \mathscr{G}r_\bullet(f)\otimes_{\mathcal{O}_Y}\mathcal{O}_{Y'}
  \longrightarrow\mathscr{G}r_\bullet(f')
\]
est bijectif (resp. surjectif).
\end{enumerate}
\end{proposition}

\noindent(i) Les hypothèses entraînent aussitôt que pour tout $y'\in Y'$,
$\rho_{y'}^*(\theta_{\psi'(y')}^\#)$ est un homomorphisme local (I, 1.6.2),
donc $w_n$ est un morphisme de préschémas (I, 2.2.1).

